\chapter*{Abstract}
\addcontentsline{toc}{chapter}{Abstract}

Carotid plaque is a marker of systemic atherosclerosis, but it is not recorded as a variable in
population cohorts: it exists only implicitly, in ultrasound images that no one has read. Any
analysis relating plaque to molecular data must therefore begin by producing the phenotype with a
model, and is bounded by the quality of that production step. This thesis sets out to identify
plasma proteins associated with subclinical carotid atherosclerosis, and examines the measurement
on which such a search depends.

A published detection model was re-run across the UK Biobank carotid ultrasound collection,
yielding detections for 179{,}611 images from 20{,}014 participants and a plaque prevalence of
43.0\,\% against the 45\,\% originally reported. Tested against the Olink proteomic panel in the
2{,}542 participants with both measurements, this phenotype showed no association surviving
correction for multiple testing: 151 of 2{,}923 proteins reached nominal significance where 146
are expected by chance, and the best classifier attained a cross-validated ROC-AUC of 0.541,
below the 0.551 of age and sex alone. Age was significantly associated in the same data,
confirming that the analysis detects association where association exists.

Because a misclassified exposure attenuates estimated associations in proportion to Youden's $J$,
the model was then evaluated outside its development cohort, on a public Argentinian dataset. Its
specificity fell from 89.2\,\% to 62.5\,\% while sensitivity declined only from 89.5\,\% to
83.3\,\%. Instance segmentation models fine-tuned on that dataset raised $J$ from 0.458 to 0.625,
and pooling two of them to 0.667, at operating points selected on validation rather than test
data. A detection control trained on the same images matched this, locating the improvement in
the adaptation rather than in the segmentation formulation.

Selecting the best-matching predicted instance irrespective of its score raises the mean Dice
coefficient from 0.655 to 0.794 and leaves no image without a correct mask: the models outline
plaques they decline to report. This was traced to the training objective, in which the
confidence term was weighted twenty-five times below the terms governing geometry. The limiting
component is therefore the image-level decision, not segmentation quality.
